\documentclass[main.tex]{subfiles}
\begin{document}



\section*{Behandlung der for-Schleife in Tests}

\subsection*{1. Behandlung der for-Schleife}
In einem Glass-Box-Test wird die for-Schleife behandelt, indem:
\begin{itemize}
    \item Jede Iteration getestet wird, um zu überprüfen, ob die Schleifenfunktionalität korrekt ist.
    \item Grenzwerte überprüft werden, um sicherzustellen, dass die Schleife ordnungsgemäß funktioniert.
    \item Leere Schleifen getestet werden, um zu sehen, wie das System darauf reagiert.
    \item Die innere Logik der Schleife getestet wird, um sicherzustellen, dass alle möglichen Ausführungspfade abgedeckt sind.
\end{itemize}

\subsection*{2. Unterschied zwischen Pfaden und Zweigen}
\begin{itemize}
    \item \textbf{Pfade:} Eine vollständige Sequenz von Anweisungen, die während der Ausführung eines Programms durchlaufen wird.
    \item \textbf{Zweige:} Entscheidungspunkte im Code, die verschiedene Ausführungsmöglichkeiten darstellen, wie z.B. die Bedingungen einer Schleife.
\end{itemize}

\end{document}